Mit $\rho_\text{Cu} = \rCuO$ und $\rho_\text{Zn} = \rZnO$ (FoTaBe S.~187) und der Annahme, dass sich die Volumina addieren, ist
\begin{align}
 \rho &= \frac{m_\text{Cu}+m_\text{Zn}}{V_\text{Cu}+V_\text{Zn}} = \rLeF \\
   &= \frac{\mCuO+\mZnO}{\frac{\mCuO}{\rCuO}+\frac{\mZnO}{\rZnO}} \\
   &= \rLeg \approx \result{\rLegP{0}{2}}
\end{align}
Die Tabelle gibt für Messing (65\,\% Cu, 35\,\% Zn) $\rho_\text{T} = \rTaO$, also um
\begin{align}
 \frac{\rho_\text{T}-\rho}{\rho} &= \dTaF = \dTa \approx \dTaP{0}{2}
\end{align}
mehr. Das Messing hat ein kleineres Volumen als Kupfer und Zink getrennt: Die Zinkatome werden ins Kristallgitter des Kupfers eingebaut und sind dort dichter gepackt als im reinen Zink. Die Volumina addieren sich deshalb nicht.
